\begin{table}[t]
\centering
\caption{Standardised Effect Sizes}
\label{tab:sde}
\begin{tabular}{lcccccc}
\toprule
Outcome & $\hat{\beta}$ & SE & SD($Y$) & SDE & SE(SDE) & Classification \\
\midrule
Log rent level & -0.0046 & 0.0089 & 0.2429 & -0.0191 & 0.0367 & Small negative \\
YoY rent growth (\%) & -2.44 & 0.77 & 4.25 & -0.5730 & 0.1810 & Large negative \\
\bottomrule
\end{tabular}
\par\vspace{0.3em}\begin{minipage}{\textwidth}\footnotesize
\textit{Notes:} \textbf{Country:} Ireland. \textbf{Research question:} Do Rent Pressure Zone designations, which cap annual rent increases at 4\%, reduce rent levels or rent growth in designated areas relative to not-yet-designated areas? \textbf{Policy mechanism:} RPZ designation prohibits landlords from raising rents on existing tenancies by more than 4\% per annum in designated Local Electoral Areas; new tenancies are subject to the same cap relative to the previous rent. \textbf{Outcome definition:} (1)~Log standardised average monthly rent (RTB/ESRI), (2)~Year-on-year percentage change in standardised average monthly rent. \textbf{Treatment:} Binary --- county designated as RPZ at staggered dates from 2016Q4 to 2021Q3. \textbf{Data:} CSO PxStat table RIQ02, quarterly, county-level, 26 counties, 2012Q1--2025Q3, $N = 1{,}430$ county-quarters. \textbf{Method:} Callaway and Sant'Anna (2021) staggered DiD; not-yet-treated control group; standard errors via multiplier bootstrap clustered at county level. \textbf{Sample:} All 26 Irish counties; restricted to 2012Q1 onwards to avoid post-crisis recovery distortion. SDE $= \hat{\beta} / \text{SD}(Y)$ where SD($Y$) is the pre-treatment standard deviation. Classification refers to magnitude, not statistical significance: Large ($|$SDE$| > 0.15$), Moderate ($0.05$--$0.15$), Small ($0.005$--$0.05$), Null ($< 0.005$).
\end{minipage}
\end{table}

